\section*{अध्याय \ref{ch:g9:trig} --- समकोण त्रिभुज में त्रिकोणमिति}

\begin{solution}{exo:g9:trig:1}
कर्ण: $\sqrt{9^2 + 12^2} = \sqrt{81 + 144} = \sqrt{225} = 15$।

दूसरी \omterm{def:g6:shapes:triangles}{समकोण} भुजा: $\sqrt{17^2 - 8^2} = \sqrt{289 - 64} = \sqrt{225} = 15$।
\end{solution}

\begin{solution}{exo:g9:trig:2}
$7^2 + 24^2 = 49 + 576 = 625 = 25^2$: हाँ, \omterm{def:g6:shapes:triangles}{समकोण} है (पाइथागोरस की
विलोम), और \omterm{def:g6:shapes:triangles}{समकोण} भुजा $25$ के सामने पड़ता है।

$5^2 + 6^2 = 61 \neq 64 = 8^2$: \omterm{def:g6:shapes:triangles}{समकोण} नहीं है।
\end{solution}

\begin{solution}{exo:g9:trig:3}
कर्ण $[EF]$ है ($D$ पर बने \omterm{def:g6:shapes:triangles}{समकोण} के सामने वाली भुजा)।
$\theta = \widehat E$ की सम्मुख भुजा $[DF]$ है, और संलग्न भुजा $[DE]$।
इसलिए
\[
\cos\theta = \frac{DE}{EF}, \qquad
\sin\theta = \frac{DF}{EF}, \qquad
\tan\theta = \frac{DF}{DE}.
\]
\end{solution}

\begin{solution}{exo:g9:trig:4}
सम्मुख \omterm{def:g6:shapes:triangles}{समकोण} भुजा: $10 \sin 28^\circ \approx 10 \times 0.469 = 4.7$।
संलग्न \omterm{def:g6:shapes:triangles}{समकोण} भुजा: $10 \cos 28^\circ \approx 10 \times 0.883 = 8.8$।
(जाँच: $4.7^2 + 8.8^2 \approx 22 + 77 \approx 100$।)
\end{solution}

\begin{solution}{exo:g9:trig:5}
$\tan\theta = \dfrac{4.5}{6} = 0.75$, इसलिए
$\theta = \tan^{-1}(0.75) \approx 37^\circ$।
\end{solution}

\begin{solution}{exo:g9:trig:6}
देखने वाले, ड्रोन के ठीक नीचे ज़मीन के बिंदु, और ड्रोन से बने \omterm{def:g6:shapes:triangles}{समकोण
त्रिभुज} में: $120$ m की ऊँचाई उन्नयन कोण की सम्मुख भुजा है और क्षैतिज
दूरी $d$ उसकी संलग्न भुजा। इसलिए
\[
\tan 32^\circ = \frac{120}{d},
\qquad
d = \frac{120}{\tan 32^\circ} \approx \frac{120}{0.625} = 192 \text{ m।}
\]
\end{solution}

\begin{solution}{exo:g9:trig:7}
चढ़ाई ($1.2$ m) कोण की सम्मुख भुजा है, और रास्ते की लंबाई $L$ कर्ण
है:
\[
\sin 5^\circ = \frac{1.2}{L},
\qquad
L = \frac{1.2}{\sin 5^\circ} \approx \frac{1.2}{0.0872} \approx 13.8
\text{ m।}
\]
इसलिए रास्ता कम से कम $13.8$ m लंबा होना चाहिए।
\end{solution}

\begin{solution}{exo:g9:trig:8}
$\cos^2\theta + \sin^2\theta = 1$ से:
$\sin^2\theta = 1 - 0.36 = 0.64$, और न्यून कोण के लिए
$\sin\theta > 0$ होता है, इसलिए $\sin\theta = 0.8$। तब
$\tan\theta = \dfrac{\sin\theta}{\cos\theta} = \dfrac{0.8}{0.6} =
\dfrac43$।
\end{solution}

\begin{solution}{exo:g9:trig:9}
\omterm{def:g6:shapes:triangles}{समकोण} भुजाएँ $1$ वाला \emph{\omterm{def:g6:shapes:triangles}{समद्विबाहु} \omterm{def:g6:shapes:triangles}{समकोण त्रिभुज}}: कर्ण
$\sqrt{1 + 1} = \sqrt2$; दोनों न्यून कोण $45^\circ$ के, और
\[
\cos 45^\circ = \frac{1}{\sqrt2} = \frac{\sqrt2}{2} = \sin 45^\circ,
\qquad
\tan 45^\circ = \frac11 = 1 .
\]

भुजा $1$ वाला \emph{\omterm{def:g6:shapes:triangles}{समबाहु त्रिभुज}}: उसकी ऊँचाई उसे दो \omterm{def:g6:shapes:triangles}{समकोण}
त्रिभुजों में बाँट देती है, जिनका कर्ण $1$, आधार $\frac12$ और ऊँचाई
$h = \sqrt{1 - \frac14} = \frac{\sqrt3}{2}$ है (पाइथागोरस)। कोण
$60^\circ$ (आधार पर) और $30^\circ$ (शिखर पर) के हैं। अनुपात पढ़ने
पर:
\[
\cos 60^\circ = \frac{1/2}{1} = \frac12, \quad
\sin 60^\circ = \frac{\sqrt3/2}{1} = \frac{\sqrt3}{2}, \quad
\cos 30^\circ = \frac{\sqrt3}{2}, \quad
\sin 30^\circ = \frac12 .
\]
\end{solution}

\begin{solution}{pb:g9:trig:1}
\textbf{1.} आँख की ऊँचाई से ऊपर की ऊँचाई $35^\circ$ की सम्मुख भुजा
है, जिसकी संलग्न भुजा $60$ m है:
$h = 60 \tan 35^\circ \approx 60 \times 0.700 = 42$ m।

\textbf{2.} $\tan\theta = 0.12$ से
$\theta = \tan^{-1}(0.12) \approx 6.8^\circ$ --- किसी सड़क के लिए
यह तीखा है, फिर भी कोण मामूली-सा। $45^\circ$ वाली सड़क पर
$\tan 45^\circ = 1$: यानी $100\,\%$ की चढ़ाई। (``सौ प्रतिशत'' खड़ी
दीवार से बहुत दूर है --- रस्सी-गाड़ी वालों की मनपसंद ग़लतफ़हमी।)

\textbf{3.} $\theta$ और $90^\circ - \theta$ न्यून कोणों वाले किसी
\omterm{def:g6:shapes:triangles}{समकोण त्रिभुज} में एक की सम्मुख भुजा दूसरे की संलग्न भुजा है, और
कर्ण दोनों में साझा। इसलिए
$\sin(90^\circ - \theta) = \frac{\text{उसकी सम्मुख भुजा}}{h} =
\frac{\theta \text{ की संलग्न भुजा}}{h} = \cos\theta$, और उसी तरह
$\cos(90^\circ - \theta) = \sin\theta$। नाम ख़ुद यही कहता है:
\emph{को}\omterm{def:g9:trig:ratios}{ज्या} असल में \emph{कोटि} की, यानी पूरक कोण की, \omterm{def:g9:trig:ratios}{ज्या} है।

\textbf{4.} ऊँचाई: $8 \sin 60^\circ = 8 \times
\frac{\sqrt3}{2} = 4\sqrt3 \approx 6.93$ m। पैर:
$8 \cos 60^\circ = 4$ m, बिलकुल सही।

\textbf{5.} $0.574^2 + 0.819^2 = 0.329 + 0.671 = 1.000$ (\omterm{def:g6:decimals:rounding}{गोलन} तक);
और $\left(\frac12\right)^2 +
\left(\frac{\sqrt3}{2}\right)^2 = \frac14 + \frac34 = 1$। जो जोड़ी
इस सर्वसमिका पर खरी न उतरे, वह या तो ग़लत पढ़ी गई है या कैलकुलेटर
ग़लत ढंग पर लगा है --- यानी मुफ़्त की ग़लती-पकड़।

\textbf{6.} $\tan 40^\circ = \dfrac hx$ और
$\tan 30^\circ = \dfrac{h}{x + 200}$।

\textbf{7.} $x = \dfrac{h}{\tan 40^\circ}$ और
$x + 200 = \dfrac{h}{\tan 30^\circ}$; घटाने पर
$200 = h\left(\frac{1}{\tan 30^\circ} -
\frac{1}{\tan 40^\circ}\right)$, और वहीं से सूत्र निकल आता है।
संख्याएँ: $\frac{1}{0.577} \approx 1.732$, $\frac{1}{0.839} \approx
1.192$, \omterm{ex:g1:subtraction:difference}{अंतर} $0.540$:
$h \approx \frac{200}{0.540} \approx 370$ m।

\textbf{8.} $x = \frac{h}{\tan 40^\circ} \approx
\frac{370}{0.839} \approx 441$ m। $A$ से जाँच:
$\frac{h}{x + 200} \approx \frac{370}{641} \approx 0.577 =
\tan 30^\circ$: मेल खाता है।

\textbf{9.} एक-ठिकाना विधि को चोटी के \emph{ठीक नीचे} वाले बिंदु
तक की क्षैतिज दूरी चाहिए --- जो किसी खाई के पार या पहाड़ के भीतर
नापी ही नहीं जा सकती। दो-ठिकाना विधि उसकी जगह वह दूरी ले लेती है
जिसे तुम ख़ुद क़दमों से नाप सकते हो: अपने दोनों ठिकानों के बीच के
$200$ m।

\textbf{10.} $h = \dfrac{200}{\sqrt3 - \frac{1}{\sqrt3}}
= \dfrac{200}{\frac{2}{\sqrt3}} = 100\sqrt3 \approx 173$ m ---
यानी एफ़िल मीनार की दूसरी मंज़िल, दो नज़रों और एक सैर से नाप ली
गई।

\textbf{11.} दृष्टि-रेखा $d$, और \omterm{def:g3:shapes:shapes}{त्रिज्याएँ} $R$ (छूने वाले बिंदु
तक, दृष्टि-रेखा पर \omterm{def:g6:lines:perp}{लंब}) तथा $R + h$ (आँख तक): पाइथागोरस से
$d^2 + R^2 = (R + h)^2 = R^2 + 2Rh + h^2$, यानी
$d^2 = 2Rh + h^2 \approx 2Rh$ ($h$, $R$ के आगे नगण्य है)। $1.80$ m
ऊँची आँखों के लिए:
$d \approx \sqrt{2 \times 6.37 \times 10^6 \times 1.8}
\approx 4\,800$ m --- समंदर का क्षितिज मुश्किल से पाँच किलोमीटर
दूर है।

\textbf{12.} $d \approx \sqrt{2 \times 6.37 \times 10^6 \times
324} \approx 64$ km। नाविकों का नियम: $3.6\sqrt{1.8} \approx
4.8$ km और $3.6\sqrt{324} = 64.8$ km --- दोनों मेल खाते हैं।

\textbf{13.} $\tan\theta = \frac{2}{60}$ से $\theta \approx
1.9^\circ$। चाँद का $0.5^\circ$ अंगूठे के क़रीब \emph{\omterm{def:g3:division:half}{चौथाई}}
हिस्से जितना है: चाँद देखने में विशाल लगता है और नाख़ून भर से ढक
जाता है --- आँख बहुत बुरा चाँदा है।

\textbf{14.} सीढ़ियाँ: $\tan^{-1}\frac{17}{29} \approx 30^\circ$
--- यानी आरामदेह सीढ़ियाँ ठीक उन्हीं ख़ास मानों वाले $30^\circ$ पर
चढ़ती हैं। $29$ cm चौड़ाई पर $60^\circ$ चाहिए तो हर क़दम
$29\tan 60^\circ = 29\sqrt3 \approx 50$ cm ऊँचा होगा: वे चढ़ने की
सीढ़ियाँ नहीं रहीं, दीवार से टिकाने वाली सीढ़ी हो गईं।

\textbf{15.} $\tan 10^\circ \approx 0.18$;
$\tan 45^\circ = 1$; $\tan 80^\circ \approx 5.7$;
$\tan 89^\circ \approx 57.3$। $\theta$ के $90^\circ$ की ओर बढ़ते
ही संलग्न भुजा सिमटकर ख़त्म होने लगती है, जबकि सम्मुख भुजा टिकी
रहती है: \omterm{def:g3:division:remainder}{भागफल} हर हद के पार चला जाता है। ठीक $90^\circ$ पर न कोई
त्रिभुज बचता है, न कोई मान --- \omterm{def:g9:trig:ratios}{स्पर्शज्या} का आलेख एक खड़े
अनंतस्पर्शी पर चढ़ जाता है, जो हाई-स्कूल खंड के ऐसे कई
अनंतस्पर्शियों में पहला है।

\end{solution}
